\section{Ethics, privacy and responsible reporting}\label{sec:ethics}

The records this paper's methods are built for describe people at a moment of victimisation. A record reveals, or implies, a person's religion or perceived religion, the fact that they were targeted, and often where they live, work or worship. Some records also concern suspects, who may be minors and who have not been convicted. This section sets out the safeguards that the analysis plan of \cref{sec:plan} commits to. It also sets out the risks that remain even when the safeguards work.

\subsection{Consent, review and agreements}\label{sec:ethics-review}

None of the lists was compiled for research, and the people recorded in them did not consent to linkage. We will seek review by an institutional review board (IRB) before any data are transferred. We will ask the IRB to decide whether a waiver of consent is justified. That decision turns on whether the safeguards below keep risk minimal and whether contacting victims would itself cause harm. The data-use agreements of \cref{sec:plan-lists} restrict use to the registered protocol. They prohibit attempts at re-identification and any linkage beyond the protocol. They require the minimum necessary fields and set retention and destruction dates; linkage keys and encodings are destroyed when the replication package is released. Data holders review outputs for disclosure risk only. In the other direction, the agreements oblige us to notify data holders of any breach. Victim-service organisations will be consulted on whether research use should be disclosed to people who report in future, for example in the privacy notice of a reporting portal. Transparency of that kind may be worth a small cost in reporting rates, but the trade-off belongs to the community and to the data holders as well as to us.

\subsection{Privacy-preserving record linkage and its limits}\label{sec:ethics-pprl}

Linkage (\cref{sec:plan-linkage}) separates roles. Each data holder encodes the identifying fields (locations and institution names; names of persons are neither requested nor needed) and sends the encodings, with non-identifying attributes, to a linkage unit inside a secure enclave. Analysts receive only capture histories and covariates. Encoding uses Bloom filters \citep{schnell2009}. Each string is split into $q$-grams, each $q$-gram is mapped by several hash functions into a bit array, and the Dice coefficient of two bit arrays approximates the $q$-gram similarity of the underlying strings. This allows approximate matching without exchanging plaintext.

Bloom-filter encodings are not anonymous. Frequency-based cryptanalysis can recover encoded values when an attacker holds a reference list whose value frequencies resemble those of the encoded data. Published attacks include a constraint-satisfaction attack \citep{kuzu2011}, an attack on the individual bit patterns of basic Bloom filters \citep{niedermeyer2014}, and fast pattern-mining attacks \citep{christen2019}. Our setting is unfavourable in a specific way. The addresses of synagogues, schools and other Jewish institutions are public, and incidents that target institutions cluster on them. A reference list therefore exists for exactly the records whose exposure would be most harmful. We will treat encodings as personal data. They will be generated with keyed hash functions whose key is held by the data holders and not by the linkage unit. They will never leave the enclave, and they will be destroyed on schedule. Where an agreement permits matching on plaintext inside the enclave, we will prefer that, because it removes the attack surface without exporting anything.

\subsection{Enclaves, aggregate release and disclosure control}\label{sec:ethics-sdc}

All record-level processing, including clerical review, takes place in a secure research environment. Nothing leaves it without output review. Only aggregates are released. Every released table is subject to small-cell suppression with a minimum cell size $c$ fixed in the data-use agreements; we propose $c=10$ \citep{hundepool2012}. Capture-history tables need care beyond single-cell suppression, because the list holders publish their own totals. NYPD open data and the H.E.A.T. Map, for example, make list sizes $n_{k+}$ approximately public for a jurisdiction. A single suppressed cell $n_{\h}$ could then be recovered by subtraction from these margins. Complementary suppression, or release only of coarsened summaries such as the frequency counts $f_j$, will be applied so that no suppressed cell is determined by released and published totals together \citep{hundepool2012}. Small-area outputs (borough by year) are released only as model-based estimates with intervals. They are never released as maps of point locations.

\paragraph{Formal differential privacy (optional)}
For any released small-area counts, we may add a formal guarantee. A randomised mechanism $\mathcal M$ is $\varepsilon_{\mathrm{DP}}$-differentially private if, for every pair of datasets $D,D'$ that differ by the addition or removal of one incident, and for every measurable set $S$,
\begin{equation}\label{eq:eth-dp}
  \Pr\{\mathcal M(D)\in S\}\le e^{\varepsilon_{\mathrm{DP}}}\,\Pr\{\mathcal M(D')\in S\}
\end{equation}
\citep{dwork2006,dwork2014}. Here $\varepsilon_{\mathrm{DP}}$ is the privacy-loss parameter, unrelated to the false-match rate $\varepsilon$ of \cref{sec:linkage}. Each incident falls in exactly one cell of the capture-history table $(n_{\h})_{\h\in\Hobs}$, so the table has $\ell_1$-sensitivity~1 under this neighbouring relation. Adding independent $\mathrm{Laplace}(0,b)$ noise with $b=1/\varepsilon_{\mathrm{DP}}$ to every cell therefore satisfies \eqref{eq:eth-dp}. Each noisy cell has standard deviation $\sqrt2\,b=\sqrt2/\varepsilon_{\mathrm{DP}}$, and any estimate computed from the noisy table inherits the guarantee by post-processing \citep{dwork2014}. The guarantee protects incidents, not people. A person who was targeted several times is protected at level $m\,\varepsilon_{\mathrm{DP}}$, where $m$ is the number of incidents involving them.

The cost of this noise is best judged against the posterior uncertainty in $N$, not against the counts. MSE estimators depend on overlap cells through ratios, so the noise is amplified when overlaps are small. For the two-list Lincoln--Petersen estimator $\Nhat=n_{1+}n_{2+}/n_{11}$, suppose $n_{11}$ is small relative to the list sizes. To first order, the noise then adds a relative standard deviation of about $\sqrt2/(\varepsilon_{\mathrm{DP}}\,n_{11})$. The sampling relative standard deviation under independent, homogeneous capture is about $\{(1-p_1)(1-p_2)/n_{11}\}^{1/2}$. The ratio of the two is
\begin{equation}\label{eq:eth-ratio}
  \frac{\mathrm{SD}_{\mathrm{DP}}}{\mathrm{SD}_{\mathrm{samp}}}\;\approx\;\frac{\sqrt2}{\varepsilon_{\mathrm{DP}}\sqrt{n_{11}(1-p_1)(1-p_2)}}.
\end{equation}
As illustrative arithmetic, not a result, take $\varepsilon_{\mathrm{DP}}=1$ and $p_1=p_2=0.3$. For $n_{11}=50$ the ratio is about $0.29$, so the total standard deviation rises by about 4\%. For $n_{11}=5$ the ratio is about $0.90$, and the total standard deviation rises by about 35\%. Differential privacy is therefore cheap where MSE is informative and expensive where MSE is already weak, which is also where suppression applies. If we use it, inference will treat the noise-free table as latent and include the Laplace likelihood in the model of \cref{sec:models}. Naive plug-in estimates can fail when a noisy overlap cell is near zero.

\subsection{Errors in either direction}\label{sec:ethics-errors}

Both overestimation and underestimation of the dark figure cause harm, and the methods can err either way. False non-matches split incidents and shrink overlaps, which pushes $\Nhat$ up. False matches inflate overlaps and push it down (\cref{thm:linkbias}). Referral pushes the Lincoln--Petersen estimate down (\cref{prop:referral}). More generally, the sign of the oriented highest-order interaction determines the direction of the bias from unmodelled dependence (\cref{prop:direction}). Scope mismatch can push them either way (\cref{ex:src-scope}).

An overestimate can cause fear out of proportion to risk and misdirect scarce security resources. When it is later corrected, it can be used to discredit legitimate monitoring and the people who report.

An underestimate erases victims' experience and makes police data look more complete than they are, which can justify inaction.

We therefore commit to three practices. First, we report the heterogeneity-robust lower bound (\cref{thm:chao}) next to every estimate. Second, we report the full sensitivity grid rather than a single preferred cell. Third, we state the direction of each known bias source alongside the result.

\subsection{Communicating uncertainty}\label{sec:ethics-communication}

Estimates of unreported hate crime will be quoted without their caveats. Our release materials will therefore follow six rules.
\begin{enumerate}[(i)]
  \item Lead with an interval and a lower bound, not a point. For example: ``at least $L$; most plausibly between $A$ and $B$''.
  \item State the scope every time: criminal incidents, the jurisdiction and the period. Do not convert to ``percentage unreported'' without the interval.
  \item Do not compare one list's published total with another's as if they counted the same thing. The gap between an NGO audit and police statistics reflects scope and unit as well as under-reporting (\cref{sec:src-harmonise}).
  \item Do not describe changes over time across definitional or reporting-system breaks as trends. Examples are the NIBRS transition of 2021 and the Audit's scope change of October 2023 (\cref{sec:src-police,sec:src-audit}).
  \item Show sensitivity visually, for example as $\Nhat$ against $\omega$ with the plausible region shaded.
  \item Provide a plain-language summary, a frequently-asked-questions document, and embargoed briefings for journalists and officials that cover what the estimate does \emph{not} show.
\end{enumerate}

\subsection{Engagement with Jewish communities}\label{sec:ethics-community}

The work concerns a community that bears the harm being measured. Before data acquisition, we will convene an advisory group. It will include Jewish communal and security organisations, synagogue and campus organisations, and victim-service providers. It will also include privacy and civil-liberties experts and researchers who study hate crime against other groups, whose data infrastructures overlap with ours. The group will be consulted on three matters: the ontology and scope variants (\cref{sec:src-harmonise}), the disclosure rules, and how results are communicated. It has no veto over findings. Funding sources and any relationship between the research team and the data holders will be disclosed in every output. Communities differ in how they relate to the police and to communal organisations. Those differences are themselves sources of capture heterogeneity, so we will seek input across denominations, languages and neighbourhoods rather than through a single channel.

\subsection{Misuse}\label{sec:ethics-misuse}

We foresee four kinds of misuse.
\begin{itemize}
  \item \emph{Targeting.} Fine-grained maps of incidents or of ``under-reporting'' could point offenders to vulnerable institutions, or to areas where offences are least likely to be recorded. We therefore release no point locations and only coarse, model-based small-area outputs.
  \item \emph{Selective quotation.} One end of an interval, or one cell of the sensitivity grid, can be quoted to support a predetermined claim. The full grid and the decision rules of \cref{sec:plan-decisions} are published so that such quotation can be checked.
  \item \emph{Definitional proxy wars.} Scope variants with and without Israel/Zionism-related content will be reported side by side. They are not a ruling on any definition of antisemitism (\cref{sec:src-harmonise}).
  \item \emph{Function creep.} A linkage infrastructure built to count victims could be repurposed to monitor people, including those engaged in lawful protest. The agreements confine linkage to the registered protocol and require keys and encodings to be destroyed. No offender-level or protester-level dataset will be created.
\end{itemize}
